\documentclass[10pt,a4paper]{article}
	
	\usepackage{amsmath,amssymb,amsthm,mathtools,amsfonts}
	\usepackage{breqn}
	\usepackage[utf8]{inputenc}
	\usepackage[english]{babel}
	\usepackage[T1]{fontenc}
	\usepackage{graphicx}
	\graphicspath{{./img/}}
	
	%\usepackage{ebgaramond}
	\usepackage{fontspec}
	\setmainfont[Numbers=OldStyle]{EB Garamond}
	\newfontfamily\mysectionfont[Numbers=OldStyle]{Klavika light}
	
	\usepackage[pdfborder={0 0 0}]{hyperref} %% Ingen firkant rundt om referencer
	\usepackage{multicol}
	\usepackage{tabularx}
	\usepackage{enumitem}% better controls with enumerating
	\usepackage{wrapfig}
		%\setlength{\wrapoverhang}{\marginparwidth}
		%\addtolength{\wrapoverhang}{\marginparsep}
		\setlength{\columnsep}{0pt}
		\setlength{\intextsep}{-10pt}
	\usepackage{caption}
	\usepackage[svgnames]{xcolor}
	\usepackage{dirtytalk}
	\usepackage[modulo]{lineno}
	\usepackage{hyperref}
	\usepackage{pdfpages}
	\usepackage{tasks}
	\newcommand\svar[1]{\newline\textcolor{red}{\textbf{Svar}\\#1}}
	\newcommand{\qm}[1]{“#1”}
	
	\definecolor{hgblue}{RGB}{10, 40, 80}
	\definecolor{hgred}{RGB}{140, 10, 10}
	\definecolor{hggrey}{RGB}{215, 215, 210}
	\definecolor{hggreen}{RGB}{145, 205, 205}
	
	\usepackage{titlesec}
	\titleformat{\subsection}{\color{hgred}\mysectionfont\large}{}{}{}
	\newcommand{\source}[1]{Source: \href{#1}{#1}}
	\usepackage{lipsum}
	\usepackage{comment}
	\usepackage{ulem}
	
	%\reversemarginpar
	\newcommand{\glos}[3]{\dotuline{#1}\marginpar{\scriptsize\textit{#3} #2}}

	\setlength{\parskip}{0.5\baselineskip}
	\setlength{\parindent}{0cm}
	\setlist[enumerate]{label=\alph*),left=20pt}
	%\setlist[enumerate]{leftmargin=\parindent}
	%\setlist[enumerate]{nosep}
	
	%Titling
	\usepackage{titling}
	\setlength{\droptitle}{-12ex}
	\pretitle{\begin{center}\mysectionfont\Huge\color{hgblue}}
	\posttitle{\end{center}}
	\preauthor{\begin{center}\mysectionfont\Large\color{hgblue}}
	\postauthor{\end{center}}
	\predate{\begin{center}\mysectionfont}
	\postdate{\end{center}}
	
	\setcounter{secnumdepth}{0}
	
	\usepackage{tikz}

\begin{document}

%\pagecolor{FloralWhite}
\title{Eksamen 3t EN}
%\date{\vspace{-5em}} % I tilfældet ingen dato
\date{Sommer, 2024}
\author{Samlede spørgsmål}
\maketitle

\begin{tikzpicture}[remember picture,overlay]
    \node[anchor=north east,yshift=-15pt,xshift=-45pt]%
       at (current page.north east)
        {
\includegraphics[scale=0.1]{HG_logo_navn 2.PNG}};
    \end{tikzpicture}

\thispagestyle{empty}

\vspace{-4em}

\begin{center}
{\color{hgblue}\rule{1\textwidth}{1pt}}
\end{center}

\subsection{Spørgsmål 1}
\begin{description}
\item[Theme:] Horror
\item[Text:] Shirley Jackson, “The Lottery” (1948) (short story) (abridged) (15819~tegn~=~3,7~ns)
\item[Instruction:] You must analyse and interpret “The Lottery” by Shirley Jackson and make a presentation (roughly 8 minutes) of the text and relate it to the theme. Apart from your own points, you must include the following points in your presentation:
\begin{itemize}
\item characterisation of the Mrs. Hutchinson 
\item identification of suspense
\end{itemize}
\end{description}

\subsection{Spørgsmål 2}
\begin{description}
\item[Theme:] Horror
\item[Text:] Edgar Allen Poe, “The Raven” (1845) (narrative poem) (abridged) (96~verselinjer~=~3,2~ns)
\item[Instruction:] You must analyse and interpret “The Raven” by Edgar Allen Poe and make a presentation (roughly 8 minutes) of the text and relate it to the theme. Apart from your own points, you must include the following points in your presentation:
\begin{itemize}
\item characterisation of the narrator
\item use of symbolism
\end{itemize}
\end{description}

\subsection{Spørgsmål 3}
\begin{description}
\item[Theme:] Horror
\item[Text:] Charles Burns, “Black Hole” (2005) (graphic novel) (excerpt)
\item[Instruction:] You must analyse and interpret the excerpt from “Black Hole” by Charles Burns and make a presentation (roughly 8 minutes) of the text and relate it to the theme. Apart from your own points, you must include the following points in your presentation:
\begin{itemize}
\item use of symbolism 
\item style of writing
\end{itemize}
\end{description}

\subsection{Spørgsmål 4}
\begin{description}
\item[Theme:] Being American
\item[Text:] Barack Obama, “A More Perfect Union” (2008) (speech) (abridged) (16138~tegn~=~3,8~ns)
\item[Instruction:] You must analyse and interpret “A More Perfect Union” by Barack Obama and make a presentation (roughly 8 minutes) of the text and relate it to the theme. Apart from your own points, you must include the following points in your presentation:
\begin{itemize}
\item style of writing
\item use of rhetorical devices
\end{itemize}
\end{description}

\subsection{Spørgsmål 5}
\begin{description}
\item[Theme:] Being American
\item[Texts:] John F. Kennedy, “Inaugural Address” (1961) (speech) and Donald J. Trump, "Inaugural Address" (2017) (speech) (15723~tegn~=~3,7~ns)
\item[Instruction:] You must analyse and interpret “Inaugural Address” by John F. Kennedy and "Inaugural Address” by Donald J. Trump, and make a presentation (roughly 8 minutes) of the texts and relate them to each other and to the theme. Apart from your own points, you must include the following points in your presentation:
\begin{itemize}
\item style of writing
\item intention and message
\end{itemize}
\end{description}

\subsection{Spørgsmål 6}
\begin{description}
\item[Theme:] The Island of Ireland
\item[Text:] Frank O'Connor, “Guests of the Nation” (1931) (short story) (abridged) (16737~tegn~=~4~ns)
\item[Instruction:] You must analyse and interpret “Gusts of the Nation” by Frank O'Connor and make a presentation (roughly 8 minutes) of the text and relate it to the theme. Apart from your own points, you must include the following points in your presentation:
\begin{itemize}
\item style of writing
\item the setting
\end{itemize}
\end{description}

\subsection{Spørgsmål 7}
\begin{description}
\item[Theme:] The Island of Ireland
\item[Text:] John Connell, “The Boat” (2017) (short story) (abridged) (16539~tegn~=~3,9ns)
\item[Instruction:] You must analyse and interpret “The Boat” by John Connell and make a presentation (roughly 8 minutes) of the text and relate it to the theme. Apart from your own points, you must include the following points in your presentation:
\begin{itemize}
\item narrative technique
\item the setting
\end{itemize}
\end{description}

\subsection{Spørgsmål 8}
\begin{description}
\item[Theme:] Postmodernism
\item[Text:] Paul Auster, “Auggie Wren's Christmas Story” (1990) (short story) (15821~tegn~=~3,8~ns)
\item[Instruction:] You must analyse and interpret “Auggie Wren's Christmas Story” by Paul Auster and make a presentation (roughly 8 minutes) of the text and relate it to the theme. Apart from your own points, you must include the following points in your presentation:
\begin{itemize}
\item characterisation of Auggie Wren
\item narrative technique
\end{itemize}
\end{description}

\subsection{Spørgsmål 9}
\begin{description}
\item[Theme:] Postmodernism
\item[Texts:] Margaret Atwood, “Happy Endings” (1983) (short story) (6866~tegn~=~1,6~ns) and Jamaica Kincaid, "Girl" (1978) (short story) (51~verse~linjer~=~1,7~ns)
\item[Instruction:] You must analyse and interpret “Happy Endings” by Margaret Atwood and "Girl" by Jamaica Kincaid and make a presentation (roughly 8 minutes) of the texts and relate them to each other and to the theme. Apart from your own points, you must include the following points in your presentation:
\begin{itemize}
\item themes and messages
\item narrative techniques
\end{itemize}
\end{description}

\subsection{Spørgsmål 10}
\begin{description}
\item[Theme:] The Shakespearean Drama
\item[Text:] William Shakespeare, “Romeo and Juliet” (drama) (excerpt) (111~verselinjer~=~3,7~ns)
\item[Instruction:] You must analyse and interpret the excerpt from act III, scene I of “Romeo and Juliet” by William Shakespeare and make a presentation (roughly 8 minutes) of the text and relate it to the theme. Apart from your own points, you must include the following points in your presentation:
\begin{itemize}
\item characterisation of Romeo 
\item language and poetic devices
\end{itemize}
\end{description}

\subsection{Spørgsmål 11}
\begin{description}
\item[Theme:] The Shakespearean Drama
\item[Texts:] William Shakespeare, “Romeo and Juliet” (drama) (95~verselinjer~=~3,1~ns) (excerpt) and "Sonnet 30" (poem) (14~verselinjer~=~0,5~ns)
\item[Instruction:] You must analyse and interpret the excerpt from act III, scene III of “Romeo and Juliet” as well as "Sonnet 30" by William Shakespeare and make a presentation (roughly 8 minutes) of the texts and relate them to each other and to the theme. Apart from your own points, you must include the following points in your presentation:
\begin{itemize}
\item characterisation of Romeo 
\item comparison of themes
\end{itemize}
\end{description}





\end{document}


\subsection{Assignment – Short Story (Fiction)}
Write an analytical essay on Haruki Murakami’s short story \qm{On Seeing the 100\% Perfect Girl One Beautiful April Morning}.

Part of your essay must focus on the effect of the short story's preoccupation with dream and reality. It should also include:
\begin{itemize}
\item A \textit{characterisation} of the main character.
\item An intepretation of the short story's \textit{theme} and \textit{message}.
\end{itemize}

Your essay must include references to the short story.

\includepdf[pages=1-]{On Seeing....pdf}